\begin{table}[htbp]
\centering
\small
\caption{Efecto TWFE de alcaldesa sobre inclusión financiera femenina (asinh)}
\label{tab:twfe}
\begin{tabular}{lcccccc}
\toprule
 & Coef & SE & p-valor & IC 95% & N & R² within \\
Outcome &  &  &  &  &  &  \\
\midrule
Contratos totales & 0.0070 & (0.0215) & 0.7466 & [-0.0353, 0.0492] & 41,905 & -0.0007 \\
Tarjetas débito & -0.0137 & (0.0213) & 0.5210 & [-0.0554, 0.0281] & 41,905 & 0.0072 \\
Tarjetas crédito & -0.0017 & (0.0172) & 0.9193 & [-0.0355, 0.0320] & 41,905 & 0.0046 \\
Créditos hipotecarios & 0.0180 & (0.0214) & 0.4002 & [-0.0239, 0.0599] & 41,905 & 0.0013 \\
Saldo total & 0.0042 & (0.0490) & 0.9312 & [-0.0918, 0.1003] & 41,905 & 0.0002 \\
\bottomrule
\end{tabular}
\vspace{2mm}
\parbox{\textwidth}{\footnotesize \textit{Nota:} Errores estándar clusterizados a nivel municipio entre paréntesis. FE de municipio y período incluidos. Control: ln(población). Escala: asinh(outcome per cápita ×10,000 mujeres adultas). *** p<0.01, ** p<0.05, * p<0.10. El coeficiente se interpreta como efecto promedio TWFE, no como ATT (ver Goodman-Bacon, 2021).}
\end{table}